\documentclass{beamer}
\usepackage{amsmath}
\usepackage{gvv}

\title{Question 12.14}
\author{AI25BTECH11040 - Vivaan Parashar}
\date{\today}

\begin{document}

\frame{\titlepage}

\begin{frame}
    \frametitle{Question: }
    4-point DFT of a real discrete-time signal $x[n]$ of length 4 is given by $X[k]$, $n = 0, 1, 2, 3$ and $k = 0, 1, 2, 3$. It is given that $X[0] = 5$, $X[1] = 1 + j1$, $X[2] = 0.5$. $X[3]$ and $x[0]$ respectively are
    \begin{multicols}{2}
        \begin{enumerate}[label=\alph*)]
            \item $1 - j$, 1.875
            \item $1 - j$, 1.500
            \item $1 + j$, 1.875
            \item $0.1 - j0.1$, 1.500
        \end{enumerate}
    \end{multicols}
\end{frame}

\begin{frame}
    \frametitle{Solution: }
    Since our sequence is of length $n=4$, let us denote $\omega_n = e^{-j\frac{2\pi}{4}} = -j$. The DFT matrix thus formed is:
    \[W_4 = \myvec{
            1 & 1 & 1 & 1\\
            1 & \omega_n & \omega_n^2 & \omega_n^3\\
            1 & \omega_n^2 & \omega_n^4 & \omega_n^6\\
            1 & \omega_n^3 & \omega_n^6 & \omega_n^9
        } = \myvec{
            1 & 1 & 1 & 1\\
            1 & -j & -1 & j\\
            1 & -1 & 1 & -1\\
            1 & j & -1 & -j
        }\]
    The inverse DFT matrix is given by:
    \[
        W_4^{-1} = \frac{1}{4} W_4^\mathrm{H} = \frac{1}{4} \myvec{
            1 & 1 & 1 & 1\\
            1 & j & -1 & -j\\
            1 & -1 & 1 & -1\\
            1 & -j & -1 & j
        }\]
\end{frame}
\begin{frame}
    (Here $A^\mathrm{H}$ denotes the Hermitian transpose of matrix A, ie the conjugate transpose of A).
    To get the DFT of our sequence, we use the formula:
    $\vec{X} = W_4\vec{x}$ or $\vec{x} = W_4^{-1}\vec{X}$ to get our sequence from our DFT.
    We know $X$ = $\myvec{5\\1+j\\0.5\\X[3]}$.
    Calculating $x[0]$:
    \[x[0] = \frac{1}{4} (5 + (1+j) + 0.5 + X[3]) = \frac{6.5 + j + X[3]}{4}\]
    Since $x[0]$ is real, the imaginary part must be zero. Thus, the imaginary part of $X[3]$ must be $-j$. Therefore, $X[3] = a - j$ where $a$ is a real number.
    Calculating $x[1]$:
    \[x[1] = \frac{1}{4} (5 + j(1+j) - 0.5 - j(a - j)) = \frac{4.5 + (1 - a)j}{4}\]
    Calculating $x[2]$:
    \[x[2] = \frac{1}{4} (5 - (1+j) + 0.5 + (a - j)) = \frac{4.5 + (a - 1) - 2j}{4}\]
\end{frame}
\begin{frame}
    Calculating $x[3]$:
    \[x[3] = \frac{1}{4} (5 - j(1+j) - 0.5 + j(a - j)) = \frac{4.5 + (a - 1)j}{4}\]
    Since $x[n]$ is real for all values of n, the imaginary parts of $x[1]$ and $x[3]$ must be zero. Thus, we have:
    \[1 - a = 0 \implies a = 1\]
    Thus, $X[3] = 1 - j$.
    Calculating $x[0]$:
    \[x[0] = \frac{6.5 + j + (1 - j)}{4} = \frac{7.5}{4} = 1.875\]
    Thus, the correct option is (a) $X[3] = 1 - j$ and $x[0] = 1.875$.
\end{frame}

\end{document}
